% ECONOMICS_PROBLEM: {"id":"C73-43","title":"Approximate stopping equilibrium with nested information","jel_code":"C73","source_id":"OP-0166"}
% FIRST_POSTED: 2026-10-10
% ATTEMPT_STATUS: solved
% ENTRY_KIND: research_attempt
% RESULT: explicit bounded two-player T=1 counterexample; no ex-ante 2^{-42}-Nash equilibrium
\documentclass[11pt]{article}
\usepackage[T1]{fontenc}
\usepackage[utf8]{inputenc}
\usepackage{lmodern}
\usepackage[margin=1in]{geometry}
\usepackage{amsmath,amssymb,amsthm,mathtools}
\usepackage{array,booktabs,hyperref}
\hypersetup{colorlinks=true,linkcolor=blue,urlcolor=blue,citecolor=blue}
\setlength{\emergencystretch}{3em}
\newtheorem{theorem}{Theorem}
\newtheorem{lemma}[theorem]{Lemma}
\theoremstyle{remark}
\newtheorem{remark}[theorem]{Remark}
\newcommand{\E}{\mathbb E}
\newcommand{\Pp}{\mathbb P}
\newcommand{\ind}{\mathbf 1}
\newcommand{\F}{\mathcal F}
\newcommand{\Hh}{\mathcal H}
\newcommand{\Maj}{\operatorname{Maj}}
\title{\texttt{C73-43}: contemporaneously nested information\ does not ensure approximate stopping equilibrium}
\author{Ji Ho Bae}
\date{10 October 2026}
\begin{document}
\maketitle
\section*{Model, process, and claim status}
\noindent\textbf{Model:} GPT Codex. \textbf{Contributor:} Ji Ho Bae.
The exact serving revision was not exposed.
The complete mathematical argument below was generated by GPT Codex,
using parallel derivation and adversarial proof checking in response to a
request for a complete hand proof of the exact catalogue statement.
No human-written mathematical edits are included. The construction and
all verification arguments are preserved in full.
This is an LLM claim of a complete counterexample; community expert review
is pending. No proof-assistant verification is claimed.
The complete original generated output is preserved at the immutable
\href{https://github.com/jbaelaw/nested-stopping-no-approximate-equilibrium/blob/6db16ed056d8f0f464dff4fcca8da2fddf3efd36/paper/Proof.tex}{manuscript TeX}
and \href{https://github.com/jbaelaw/nested-stopping-no-approximate-equilibrium/blob/6db16ed056d8f0f464dff4fcca8da2fddf3efd36/paper/Proof.pdf}{manuscript PDF}.


\begin{abstract}
We construct a bounded two-player stopping game with dates $0,1$ on a
standard Borel common-prior probability space. Its observation filtrations
are contemporaneously nested, but it has no measurable randomized ex-ante
$2^{-42}$-Nash equilibrium. The counterexample includes every simultaneous
stopper set and the no-stop outcome. The proof is self-contained: a
quantitative obstruction to a measurable noisy-majority recursion yields
a binary Bayesian game without approximate equilibrium; a hidden nature
flag creates a common indifference row; and that row permits an exact
realization as a stopping game. No nonexistence theorem for another game
is invoked.
\end{abstract}

\section{The full target and the result}

Problem \texttt{C73-43} asks whether every bounded finite-horizon stopping
game on an arbitrary common-prior probability space, with
\[
 \F_{n,t}\subseteq\cdots\subseteq\F_{1,t}\qquad(0\le t\le T),
\]
admits an ex-ante $\varepsilon$-Nash equilibrium for every
$\varepsilon>0$. Each player has independent private uniform seeds,
independent of the underlying state. A strategy is a randomized stopping
time in $\{0,\ldots,T,\infty\}$. The first stopping date and the full
simultaneous stopper set determine a bounded measurable payoff vector;
a separate bounded vector is paid if nobody stops. Public survival is
observed. The payoff variables in the displayed problem are arbitrary
bounded measurable state functions; no measurability with respect to an
individual player's private filtration is imposed. One-step revelation
of the richer information to the less-informed player is also not
imposed.

\begin{theorem}\label{thm:main}
There is a game in this entire specified family with two players and
$T=1$ which has no measurable randomized ex-ante $2^{-42}$-Nash
equilibrium. Its common-prior space is standard Borel, and its payoff
bounds can be chosen as $3/2$ for Player~1 and $1$ for Player~2.
Consequently, the universal existence assertion in
\texttt{C73-43} is false.
\end{theorem}

The contemporaneously nested stopping question was posed in
Section~4 of the 2024 version of Jacobovic--Solan
\cite{JS}. Simon--Tomkowicz's work on Bayesian nonexistence
\cite{ST} supplies relevant context, but its nonexistence theorem is not
a premise of the argument here. Prior \texttt{C73-43} finite-compression
attempts retained an unresolved status: they proved equilibrium transfer
under payoff and prediction-error control, and established that control
under additional product-domination assumptions. Those conditional
results did not settle the universal target. This paper answers that
target negatively by a complete explicit construction. No novelty claim
is made.

The earlier conditional and subclass results are preserved. In the direct
type encoding used here, the joint signal prior is not dominated by the
product of the marginal signal laws. On phase
$R$, of probability $1/2$, Player~1's tree signal is exactly the tree
coordinate of Player~2's signal. Under the product of the marginal laws,
those two tree coordinates are independent with the nonatomic law $m$,
so equality has probability zero. The non-atomicity follows already from
any one uniform tree coordinate. This observation about the direct type
encoding makes no assertion about alternative primitive-coordinate
representations or about the scope of the earlier coordinate-model
theorem.

\section{A quantitative obstruction on a labeled tree}

Let $W=\{1,2,3\}^{<\mathbb N}$ be the countable set of finite words,
including the empty word $e$. Define
\[
 X=(\{0,1\}\times[0,1]^3)^W.
\]
At vertex $w$, write its four coordinates as
$(b_w,t_{w,1},t_{w,2},t_{w,3})$. Equip $X$ with its countable product
Borel sigma field and the product probability $m$, with each $b_w$ a
fair bit and every $t_{w,i}$ uniform on $[0,1]$, all independent.
This is a standard Borel probability space. For $i\in\{1,2,3\}$,
define the subtree map $T_i$ by
\[
 b_w(T_i x)=b_{iw}(x),\qquad
 t_{w,j}(T_i x)=t_{iw,j}(x).
\]
The three random trees $T_1x,T_2x,T_3x$ are independent with law $m$,
and are independent of all root coordinates. Set
$s(x)=(-1)^{b_e(x)}$. For three signs,
\[
 \Maj(z_1,z_2,z_3)=\operatorname{sign}(z_1+z_2+z_3)
 =\frac{z_1+z_2+z_3-z_1z_2z_3}{2},
 \qquad z_i\in\{-1,1\}.
\]

\begin{lemma}\label{lem:tree}
There is no measurable $z:X\to\{-1,1\}$ for which
\[
 \delta:=m\{z(x)\ne s(x)\Maj(z(T_1x),z(T_2x),z(T_3x))\}
 \le\frac1{4096}.
\]
\end{lemma}
\begin{proof}
Put
\[
 z^*(x)=s(x)\Maj(z(T_1x),z(T_2x),z(T_3x)),\qquad
 \mu=\E_m z.
\]
The root sign is fair and independent of the child trees, so
$\E_m z^*=0$. Since $z,z^*$ are signs,
\begin{equation}\label{eq:treeerror}
 |\mu|\le\E_m|z-z^*|=2\delta,
 \qquad \|z-z^*\|_2=2\sqrt\delta.
\end{equation}
Let $\Hh_n$ be generated by the coordinates at vertices of length less
than $n$; $\Hh_0$ is trivial. Put
\[
 a_n=\E_m[z\mid\Hh_n],\qquad v_n=\E_m a_n^2,
 \qquad a_{n,i}(x)=a_n(T_i x).
\]
Conditionally on $\Hh_{n+1}$, the unexplored portions of the three
child trees are independent. Thus
\begin{equation}\label{eq:treeconditional}
 \E_m[z^*\mid\Hh_{n+1}]
 =\frac{s(x)}2
  (a_{n,1}+a_{n,2}+a_{n,3}-a_{n,1}a_{n,2}a_{n,3}).
\end{equation}
The variables $a_{n,1},a_{n,2},a_{n,3}$ are independent and identically
distributed, each with mean $\mu$ and second moment $v_n$.
Expanding the square in \eqref{eq:treeconditional} gives
\begin{align}
 \left\|\E_m[z^*\mid\Hh_{n+1}]\right\|_2^2
 &=\frac14(3v_n+6\mu^2-6v_n\mu^2+v_n^3)\notag\\
 &=\frac{3v_n+v_n^3}{4}+\frac32\mu^2(1-v_n).
 \label{eq:treevariance}
\end{align}
Conditional expectation contracts the $L^2$ norm. The triangle
inequality and \eqref{eq:treeerror}--\eqref{eq:treevariance} imply
\begin{equation}\label{eq:treerecursion}
 \sqrt{v_{n+1}}
 \le\sqrt{\frac{3v_n+v_n^3}{4}+\frac32\mu^2(1-v_n)}
       +2\sqrt\delta.
\end{equation}
We prove $v_n\le1/4$ for every $n$. Initially,
$v_0=\mu^2\le4\delta^2<1/4$. If $v_n\le1/4$, then
\[
 \frac{3v_n+v_n^3}{4}+\frac32\mu^2(1-v_n)
 \le\frac{49}{256}+6\delta^2
 <\left(\frac{15}{32}\right)^2,
 \qquad 2\sqrt\delta\le\frac1{32}.
\]
Hence \eqref{eq:treerecursion} gives
$\sqrt{v_{n+1}}<1/2$, completing the induction.

The sigma field generated by all $\Hh_n$ is the full product sigma
field. Finite-coordinate functions are dense in $L^2(m)$, and
conditional expectations fix each such function for all sufficiently
large $n$. Since conditional expectation is a contraction, this density
implies $a_n\to z$ in $L^2$. Therefore
$v_n\to\E_m z^2=1$, contradicting $v_n\le1/4$.
\end{proof}

\section{A two-player binary Bayesian game}

Let the state space $\Omega_0$ be the disjoint union
\[
 G=X\times\{1,2,3\},\qquad
 R=X\times\{1,2,3\}\times[0,1].
\]
Give phases $G,R$ probability $1/2$ each. Conditionally on $G$,
$(x,i)$ has law $m$ times the uniform law on $\{1,2,3\}$.
Conditionally on $R$, $(y,i,t)$ has the corresponding product law
with Lebesgue measure for $t$.

Both players have actions $\{0,1\}$. Player~1 observes $x$ in
state $(G,x,i)$, and $y$ in state $(R,y,i,t)$. Player~2 observes
\begin{equation}\label{eq:signals}
 (i,T_i x,t_{e,i}(x))\quad\hbox{in phase }G,
 \qquad (i,y,t)\quad\hbox{in phase }R.
\end{equation}
The two conditional laws of Player~2's signal agree: each is the
uniform law on $i$, times $m$, times Lebesgue measure. Player~1's two
conditional signal laws also agree. Consequently, conditional on either
player's signal, the two phases each have probability $1/2$.

The state-dependent payoffs are
\begin{equation}\label{eq:basepayoffs}
\begin{array}{c|cc}
 \text{state}&u_1(\omega,a_1,a_2)&u_2(\omega,a_1,a_2)\\ \hline
 (G,x,i)&a_1s(x)(3a_2-3/2)&0\\
 (R,y,i,t)&0&a_2(a_1-t).
\end{array}
\end{equation}
These functions are bounded and Borel measurable. A measurable mixed
strategy of Player~1 is a measurable $p:X\to[0,1]$; that of
Player~2 consists of measurable
$q_i:X\times[0,1]\to[0,1]$. Each value is the probability of
action $1$. Define
\[
 Q_i(x)=q_i(T_i x,t_{e,i}(x)),\qquad
 D(x)=s(x)\left(\sum_{i=1}^3Q_i(x)-\frac32\right).
\]
The maximal ex-ante unilateral gains $R_1,R_2$ are
\begin{equation}\label{eq:regret1}
 R_1=\frac12\E_m
   \left[|D(x)|\,|p(x)-\ind_{\{D(x)>0\}}|\right],
\end{equation}
\begin{equation}\label{eq:regret2}
 R_2=\frac16\sum_{i=1}^3\E_m\int_0^1
   |p(y)-t|\,|q_i(y,t)-\ind_{\{p(y)>t\}}|\,dt.
\end{equation}
Indeed, integrating the sampled index in phase $G$ gives Player~1's
payoff difference $D/2$. Phase $R$ gives Player~2's payoff difference
$(p(y)-t)/2$ at each signal. Choosing the action with positive payoff
difference is a measurable deviation, and pointwise maximization gives
\eqref{eq:regret1} and \eqref{eq:regret2}. A zero difference causes no
difficulty: either action is optimal there.

\begin{lemma}\label{lem:base}
The game \eqref{eq:basepayoffs} has no measurable mixed
$2^{-40}$-Nash equilibrium.
\end{lemma}
\begin{proof}
Suppose $R_1,R_2\le\varepsilon$ for some $0<\varepsilon\le1$.
Put
\[
 h_i(x)=\ind_{\{p(T_i x)>t_{e,i}(x)\}},\qquad
 E(x)=\sum_{i=1}^3|Q_i(x)-h_i(x)|.
\]
For fixed $y$, the interval
$\{t\in[0,1]:|p(y)-t|<\eta\}$ has length at most $2\eta$.
Outside that interval, the unweighted relay error is bounded by its
weighted error divided by $\eta$. Since
$(T_i x,t_{e,i}(x))$ has law $m\otimes\mathrm{Leb}$,
\eqref{eq:regret2} gives
\[
 \E_m E\le6\eta+\frac{6\varepsilon}{\eta}.
\]
Take $\eta=\sqrt\varepsilon$. Markov's inequality then gives
\begin{equation}\label{eq:relayerror}
 \E_m E\le12\sqrt\varepsilon,
 \qquad m\{E>1/4\}\le48\sqrt\varepsilon.
\end{equation}
On $\{E\le1/4\}$, $\sum Q_i$ is within $1/4$ of the integer
$\sum h_i$. Its distance from $3/2$ is at least $1/4$, so
$|D|\ge1/4$ on this set. Let
$c(x)=\ind_{\{p(x)>1/2\}}$. Rounding to the closest endpoint gives
$|p-c|\le|p-a|$ for every $a\in\{0,1\}$, as well as
$|p-c|\le1/2$. Therefore
\begin{align}
 \alpha:=\E_m|p-c|
 &\le\frac12m\{E>1/4\}
       +4\E_m[|D|\,|p-\ind_{\{D>0\}}|]\notag\\
 &\le24\sqrt\varepsilon+8\varepsilon.
 \label{eq:rounding}
\end{align}
Each root threshold $t_{e,i}$ is uniform and independent of $T_i x$.
For fixed $y$, $\ind_{\{p(y)>t\}}$ disagrees with $c(y)$ on a set
of threshold values of length exactly $|p(y)-c(y)|$. Hence
\begin{equation}\label{eq:childrounding}
 m\{h_i(x)\ne c(T_i x)\}=\alpha.
\end{equation}
Define
\[
 B(x)=\ind_{\{s(x)(\sum_i c(T_i x)-3/2)>0\}}.
\]
If $E\le1/4$ and all $h_i=c(T_i x)$, then
$\ind_{\{D>0\}}=B$ and $|D|\ge1/4$. If additionally $c\ne B$,
then $|p-B|\ge1/2$. The integrand in \eqref{eq:regret1} before its
factor $1/2$ is thus at least $1/8$. The union bound and
\eqref{eq:relayerror}--\eqref{eq:childrounding} now imply
\begin{align}
 m\{c\ne B\}
 &\le48\sqrt\varepsilon+3\alpha+16\varepsilon\notag\\
 &\le120\sqrt\varepsilon+40\varepsilon
 \le160\sqrt\varepsilon.
 \label{eq:colourerror}
\end{align}
For $z=2c-1$, the identity $c=B$ is precisely
\[
 z(x)=s(x)\Maj(z(T_1x),z(T_2x),z(T_3x)).
\]
Taking $\varepsilon=2^{-40}$ in \eqref{eq:colourerror} yields
\[
 \delta\le160\,2^{-20}=\frac5{32768}<\frac1{4096},
\]
contradicting Lemma~\ref{lem:tree}.
\end{proof}

\section{Creating a common indifference row}

Independently adjoin a fair flag $f\in\{o,d\}$ to $\Omega_0$,
creating the state space $\Omega=\Omega_0\times\{o,d\}$.
Player~1 observes its former signal together with this flag. Player~2
observes exactly \eqref{eq:signals} and does not observe the flag.
Write Player~1's two mixed strategies as $p(x)$ on flag $o$ and
$r(x)$ on flag $d$. Player~2 still uses $q_i(y,t)$.

Define the new payoffs $F_1,F_2$ by
\begin{equation}\label{eq:flagpayoffs}
\begin{array}{c|cc}
 \text{state}&F_1(\omega,a_1,a_2)&F_2(\omega,a_1,a_2)\\ \hline
 o,G,x,i&a_1s(x)(3a_2-3/2)&0\\
 o,R,y,i,t&0&a_1a_2\\
 d,G,x,i&a_1&0\\
 d,R,y,i,t&a_1&-ta_1a_2.
\end{array}
\end{equation}
In every state,
\begin{equation}\label{eq:rowzero}
 F_j(\omega,0,0)=F_j(\omega,0,1)=0,\qquad j=1,2.
\end{equation}
The bounds are $|F_1|\le3/2$ and $|F_2|\le1$.

\begin{lemma}\label{lem:flag}
The game \eqref{eq:flagpayoffs} has no measurable mixed
$2^{-42}$-Nash equilibrium.
\end{lemma}
\begin{proof}
Suppose a profile is a $\rho$-equilibrium. Player~1 may deviate only
on flag $d$, choosing action $1$ there. The gain is
$(1/2)\E_m(1-r)$, so
\begin{equation}\label{eq:dummyerror}
 \E_m(1-r)\le2\rho.
\end{equation}
Player~1 may instead change only its strategy on flag $o$. Each such
gain is one half its gain in the base game
\eqref{eq:basepayoffs}. The projected profile $(p,q)$ consequently
has base-game regret $R_1\le2\rho$.

Let $q'$ be any measurable deviation of Player~2. Write $g$ for its
gain in \eqref{eq:basepayoffs} and $\widehat g$ for its gain in
\eqref{eq:flagpayoffs}. With $\E_{y,i,t}$ denoting expectation under
$m$, uniform $i$, and Lebesgue $t$, these gains satisfy
\[
 \widehat g=\frac14\E_{y,i,t}
       (q'_i-q_i)(p(y)-t r(y)),\qquad
 \frac12g=\frac14\E_{y,i,t}
       (q'_i-q_i)(p(y)-t).
\]
Since $|q'_i-q_i|\le1$ and $0\le t\le1$,
\eqref{eq:dummyerror} gives
\[
 \left|\widehat g-\frac12g\right|
 \le\frac14\E_m(1-r)\le\frac\rho2.
\]
The equilibrium bound $\widehat g\le\rho$ implies $g\le3\rho$.
Thus $R_2\le3\rho$, and $(p,q)$ is in particular a
$4\rho$-equilibrium of the base game. At $\rho=2^{-42}$ this
contradicts Lemma~\ref{lem:base}.
\end{proof}

\section{Exact stopping-game realization}

Work on the state space $\Omega$ of \eqref{eq:flagpayoffs} with its
Borel sigma field and the specified common prior. Let $A$ be generated
by Player~1's signal, comprising the flag and its tree signal. Let $B$
be generated by Player~2's signal \eqref{eq:signals}. Set
\begin{equation}\label{eq:filtrations}
 \F_{1,0}=A,\quad \F_{1,1}=A\vee B,\qquad
 \F_{2,0}=\{\varnothing,\Omega\},\quad \F_{2,1}=B.
\end{equation}
These are increasing filtrations, and
$\F_{2,t}\subseteq\F_{1,t}$ for both dates. If later dates are
desired, keep the filtrations constant; the counterexample itself uses
$T=1$.

Write $F(a,b)$ for the vector
$(F_1(\omega,a,b),F_2(\omega,a,b))$. Assign the first-stopping
payoff vectors $X_{J,t}$ and no-stop vector $Y$ by
\begin{equation}\label{eq:stoppingpayoffs}
\begin{array}{c|ccc}
 t&X_{\{1\},t}&X_{\{2\},t}&X_{\{1,2\},t}\\ \hline
 0&F(0,0)=0&F(1,1)&F(0,1)=0\\
 1&F(1,0)&F(1,1)&F(1,1)
\end{array}
\qquad Y=F(1,0).
\end{equation}
For player $i$, take the $i$th component as $X_{i,J,t}$ or $Y_i$.
Every entry is a bounded measurable state function, precisely as
required by \texttt{C73-43}.

To include public survival without introducing spurious dependence on
another player's seed, first express a strategy as its prescriptions on
the two possible decision dates. Its date-$0$ prescription is a function
of its initial signal and private seed. There is only one surviving public
history at date $1$: both players continued at date $0$. The prescription
on that history is a function of its date-$1$ private information and seed;
its own earlier action and private memory are also determined by these
inputs. Extend this survivor-history prescription counterfactually to all
states and seeds, including paths on which the game has already ended.
Define the resulting potential stopping time $\tau_i$ as $0$ if its
date-$0$ prescription stops, otherwise $1$ if the extended date-$1$
prescription stops, and otherwise $\infty$. This potential stopping time
is adapted to the player's own filtration and seed, and agrees with its
actual play up to termination. Thus the potential stopping times induce
exactly the same realized payoff. If strategies are given directly as the
randomized stopping times in the problem statement, use those stopping
times themselves.

For any such potential randomized stopping profile, define
\begin{equation}\label{eq:actionsfromstops}
 a=\ind_{\{\tau_1>0\}},\qquad b=\ind_{\{\tau_2\le1\}}.
\end{equation}
The realized stopping payoff is exactly $F(a,b)$, pointwise in both the
state and the seeds. The exhaustive verification is as follows.
If $\tau_1=0$, Player~1 belongs to the date-$0$ stopper set, and
every outcome pays zero, which is row $a=0$ in \eqref{eq:rowzero}.
If $\tau_1>0$ and $\tau_2=0$, Player~2 alone stops first at date
$0$, giving $F(1,1)$. If both stopping times exceed $0$, a date-$1$
stop by Player~2, alone or tied, gives $F(1,1)$; a date-$1$ stop by
Player~1 alone gives $F(1,0)$; and no stop gives $F(1,0)$. These
are all possible cases, including simultaneous stops at both dates.

For each fixed private seed, the event $\{\tau_1=0\}$ is
$\F_{1,0}$-measurable, and $\{\tau_2\le1\}$ is
$\F_{2,1}$-measurable. Thus $a$ depends only on Player~1's initial
signal and its own seed; $b$ depends only on Player~2's signal and its
own seed. Integrating the seeds produces measurable mixed strategies
of \eqref{eq:flagpayoffs}. Their actions are independent conditional
on the state because the seeds are independent. This is exactly the
Bayesian game's mixed extension, even though Player~1 can acquire
additional information at date $1$: that information can change only
whether it stops at $1$ or never stops, and the payoff assignment
already gives $F(1,b)$ in either case.

The counterfactual extension also shows why public survival preserves
this reduction. In any path where Player~1 stopped at date $0$, its
binary action is $a=0$, so its payoff is independent of Player~2's
counterfactual action. If Player~2 stopped at date $0$ while Player~1
continued, then $a=b=1$ already determines the payoff, regardless of
Player~1's counterfactual date-$1$ prescription. Conditioning on survival
changes beliefs, but its incentive effect is already present in the
ex-ante payoff expression $F(a,b)$. There is no additional surviving
public branch or omitted signaling action.

Conversely, every measurable mixed strategy of
\eqref{eq:flagpayoffs} has an admissible stopping implementation.
Using independent uniform seeds, Player~1 stops at $0$ exactly when
its binary action is $0$, otherwise never stopping. Player~2 stops at
$1$ exactly when its binary action is $1$, otherwise never stopping.
These are randomized stopping times for \eqref{eq:filtrations}.
The same implementation works for each unilateral Bayesian deviation,
and \eqref{eq:stoppingpayoffs} preserves its expected gain exactly.
Therefore a stopping $\rho$-equilibrium would induce a Bayesian
$\rho$-equilibrium in \eqref{eq:flagpayoffs}. Lemma~\ref{lem:flag}
excludes this at $\rho=2^{-42}$, proving Theorem~\ref{thm:main}.

\section{Scope, measurability, and normalization}

The obstruction concerns the ex-ante equilibrium notion in the target:
every measurable unilateral deviation must have gain at most the fixed
$\varepsilon$. It already excludes this integrated notion, and hence
also any stronger equilibrium notion that implies it. It treats all
admissible randomized stopping strategies, not a restricted strategy
class.

The construction satisfies arbitrary bounded state-dependent stopping
payoffs and contemporaneously nested private information. It does not
assert the same counterexample under an additional assumption that all
payoff variables are measurable with respect to the less-informed
player's private filtration, or that both initial private sigma fields
are trivial. Neither restriction is present in the specified target.

For payoff range $[0,1]$, replace every outcome payoff vector $V$ in
the stopping game by
\[
 \frac13V+(1/2,1/2).
\]
The stated bounds ensure that both components then lie in $[0,1]$.
This divides every gain by three and preserves all comparisons, so the
normalized game has no $(2^{-42}/3)$-equilibrium. The zero row becomes
a common constant row, which preserves the exact realization.

All conclusions remain valid if the observation sigma fields are
completed with respect to the common prior. Measurable strategy
functions then have versions measurable under the underlying standard
Borel signal laws. Changes on null signal sets do not alter the payoff
integrals or the relay estimates: every subtree map $T_i$ preserves
$m$, and every relay signal has the stated product distribution.

\begin{thebibliography}{9}
\bibitem{JS}
R.~Jacobovic and E.~Solan,
\emph{Bayesian Games with Nested Information},
arXiv:2402.14450v1 (2024), Section~4, especially
``Multi-stage Bayesian games'' and
``Stopping games with asymmetric information.''
\url{https://arxiv.org/html/2402.14450v1}.

\bibitem{ST}
R.~S.~Simon and G.~Tomkowicz,
\emph{A Bayesian Game without $\varepsilon$-equilibria},
arXiv:1702.02090; Israel Journal of Mathematics \textbf{227}
(2018), 215--231.
\url{https://arxiv.org/abs/1702.02090}.
\end{thebibliography}
\end{document}
